\documentclass{amsart}
\usepackage[utf8]{inputenc}
    
\usepackage{titlesec}
\titleformat{\section}{\normalfont\bf\Large}{\Roman{section}. }{0em}{}
\titleformat{\subsection}{\normalfont\bf\large }{\relax}{0em}{}
\titleformat{\subsubsection}{\normalfont\bf}{\relax}{0em}{}

\usepackage{enumitem}
\usepackage{quoting}

\usepackage{xcolor}
\long\def\gray#1{{\color{gray}{#1}}}\let\grey=\gray
\long\def\blue#1{{\color{blue}{#1}}}
\long\def\red#1{{\color{red}{#1}}}

\usepackage{wrapfig}
\usepackage{lpic}

\usepackage{manyfoot}
\SetFootnoteHook{\hspace{1.5em}}
\DeclareNewFootnote{A}
\SetFootnoteHook{\hspace{1.5em}}
\DeclareNewFootnote{E}[alph]
\usepackage[hyperfootnotes=false,
    colorlinks,
    linkcolor={blue!80!black},
    citecolor={blue!50!black},
    urlcolor={blue!80!black}]{hyperref}
\def\foothref#1#2{\href{#1}{ #2}\footnoteA{\url{#1}}}

\newcommand{\skipthis}[2][\dots]{\red{[#1]}\gray{#2}}
\newtheorem*{theorem}{Theorem}
\newtheorem*{question}{Question}
\newtheorem*{problem}{Problem}
\newtheorem{exercise}{Exercise}

\let\dmo=\DeclareMathOperator

\let\tmp\epsilon
\let\epsilon\varepsilon
\let\varepsilon\tmp


\begin{document}
\centerline{\bf\Large From Ramsey’s Theorem to the Erdős–Hajnal
  Conjecture I%
  \footnoteE{English translation by Rostislav Matveev of the article
    \textit{Du Th\'eor\`eme de Ramsey \`a la Conjecture d’Erdős-Hajnal (1)},
    by  Patrice Ossona de Mendez, (2017),
    \url{https://images-archive.math.cnrs.fr/%
      Du-Theoreme-de-Ramsey-a-la-Conjecture-d-Erdős-Hajnal-1.htm}.
  }}%
\medskip

\centerline{by Patrice Ossona de Mendez}
\red{[Figure out how to deal with url]}
\begin{center}
  \begin{lpic}[draft,b(12ex),t(5ex),clean]{Pix/ramsey(0.5)}
     \lbl[t]{70,-5;\parbox{0.9\textwidth}{
        \begin{center}
          \sf Ordnung muss sein! (Order is inevitable.) This saying,
          attributed to Theodore Motzkin, aptly summarizes Ramsey
          theory, which aims to solve problems of the form: ``From what
          size on can a structure no longer avoid having a certain
          property?''
        \end{center}
      }}
  \end{lpic}
\end{center}

\section{A first example}

A classic recreational problem is the “friends problem”:
\begin{quoting}[leftmargin=5mm,rightmargin=5mm]\sf
  When six people meet, we are sure that either three of them know one
  another, or three of them do not know one another.
\end{quoting}
The proof of this fact is a bit tedious, because one has to consider
different cases. Nevertheless, this annoying complication will open
the way to a more general and more interesting point of view. Let us
look at it more closely:

So we have six people. Let us call Dominique%
\footnoteA{Note that this is an epicene name, as gender is of no
  relevance here.} %
one of the people in the group.

If Dominique knows at least three people in the group, two cases can occur:
\begin{itemize}
\item either two of them know each other (and together with Dominique,
  that makes at least three people who know one another);
\item or they do not know each other (and that makes three people who
  do not know one another).
\end{itemize}
We see that in both cases we reach the desired conclusion. So we only
have to consider the case where Dominique knows at most two people in
the group. Again, two cases can occur among the people whom Dominique
does not know:
\begin{itemize}
\item either two of them do not know each other (and together with
  Dominique, that makes three people who do not know one another);
\item or they all know one another (and that makes at least three
  people who know one another).
\end{itemize}
In this case as well, we reach the desired conclusion.



\end{document}
